\section{Conclusion}
\label{sec:conclusion}

We have presented \sys{}, a formally specified multi-stage platform for automated
academic \LaTeX{} document generation, and have demonstrated that formal specification
at the design level translates into measurable, reproducible runtime guarantees.

The central architectural insight is that document-generation quality is not a single
property but a product of four independently measurable sub-properties — compilability,
citation integrity, visual coherence, and language consistency — each of which is
maximised by a distinct formal component of the Unified System Model $\mathcal{P}$.
The \cgrag{} gate $\Gamma$ is responsible for citation integrity; the \tvg{} grammar $G$
for visual coherence and compilability; the Pipeline Algebra $\Pi$ and
\texttt{concurrentMap} for order preservation; and the Language Partition $\mathcal{L}_{25}$
for language consistency.

The two main theorems provide closed-form performance guarantees that are
directly actionable for system designers:
Theorem~\ref{thm:repair} bounds the figure compilation failure probability as
$(1-p)^{K+1}$, enabling principled selection of the attempt budget $K$;
Theorem~\ref{thm:concurrent} proves order preservation of concurrent execution
and provides the Amdahl speedup formula for capacity planning.

Evaluation on 120 live user sessions confirms that the full pipeline achieves
$\bar{Q}(s_5)=0.97$ — a $29\%$ relative improvement over reference-free generation
and $10\%$ over figure-free generation — and that all five system invariants hold
with 100\% satisfaction (I1--I4) and 98.3\% satisfaction (I5, limited by transient
network failures).
The \tvg{} repair cycle achieves $99.2\%$ cumulative figure compilation success
within four attempts, within $0.2$ percentage points of the theoretical prediction.

The formal model $\mathcal{P}$ is released as a citable specification in this report.
We invite the research community to extend its components — in particular the
lattice-theoretic formulation of $\Gamma$ and the MDP framing of Proposition~\ref{prop:mdp}
— as a foundation for further work on principled LLM-orchestration systems for
scholarly document production.

\paragraph{Acknowledgements.}
The author thanks the user community of \sys{} for generating the session corpus used
in this evaluation, and the Astana IT University HPC cluster for providing the
computational resources for the concurrency experiments.
